\section{结论展望}

\begin{frame}{研究结论}
  \begin{enumerate}
    \item 绿色金融发展对碳排放强度具有显著抑制作用，基准结果为 \hl{-2.363*}。
    \item 工具变量与多种稳健性检验均支持该结论，说明结论具有较强可靠性。
    \item 绿色金融减排效应存在明显区域异质性，\hlb{西部显著、东中部不显著}。
    \item 因此，绿色金融政策设计应从统一扩张转向 \hlt{区域差异化配置}。
  \end{enumerate}
\end{frame}

\begin{frame}{局限与未来研究方向}
  \begin{columns}[T,onlytextwidth]
    \column{0.48\textwidth}
    \begin{block}{当前局限}
      \begin{itemize}
        \item 绿色金融使用综合指数，难区分不同工具效应
        \item 样本停留在省级层面，时间跨度有限
        \item 因果识别仍有进一步加强空间
      \end{itemize}
    \end{block}

    \column{0.48\textwidth}
    \begin{exampleblock}{未来方向}
      \begin{itemize}
        \item 下沉到地级市或企业层面
        \item 细分绿色信贷、债券、保险等工具
        \item 引入双重差分等更强识别策略
      \end{itemize}
    \end{exampleblock}
  \end{columns}

  \vspace{0.6em}
  \begin{alertblock}{一句话总结}
    本文证明了绿色金融不仅是政策倡议，更是能够被实证识别的碳减排工具，但要真正发挥作用，关键在于 \hl{分地区、分工具、分场景优化配置}。
  \end{alertblock}
\end{frame}
